%% LyX 2.4.4 created this file.  For more info, see https://www.lyx.org/.
%% Do not edit unless you really know what you are doing.
\documentclass[oneside,english]{book}
\usepackage[T1]{fontenc}
\usepackage[latin9]{inputenc}
\setcounter{secnumdepth}{3}
\setcounter{tocdepth}{3}
\usepackage{amsthm}
\usepackage{amssymb}
\usepackage{geometry}
\geometry{verbose,tmargin=1in,bmargin=1in,lmargin=1.5in,rmargin=1.5in}
\PassOptionsToPackage{normalem}{ulem}
\usepackage{ulem}

\makeatletter
%%%%%%%%%%%%%%%%%%%%%%%%%%%%%% Textclass specific LaTeX commands.
\theoremstyle{definition}
\newtheorem{defn}{\protect\definitionname}

%%%%%%%%%%%%%%%%%%%%%%%%%%%%%% User specified LaTeX commands.
\usepackage{hyperref}
\usepackage[usenames,dvipsnames]{xcolor} %adds additional color options
\hypersetup{colorlinks=true, urlcolor=Blue}%Blue

\usepackage{multicol}

\usepackage{tikz}
\usepackage{tikz-3dplot}
\usetikzlibrary{calc,arrows,shapes,backgrounds}

\usetikzlibrary{patterns}
\usetikzlibrary{plotmarks}
\usepackage{pgfplots}
\pgfplotsset{compat=newest}

\usepackage{calc}
\usepackage[nomessages]{fp}

\usepackage{datetime}
%\usepackage{subcaption}

%\usepackage{amsthm}
% Added by lyx2lyx
\setlength{\parskip}{\smallskipamount}
\setlength{\parindent}{0pt}

\makeatother

\usepackage{babel}
\providecommand{\definitionname}{Definition}

\begin{document}

\chapter{Homework Hints}

\section*{Chapter \ref{chap:Ideals}}
\begin{itemize}
\item There was a typo in problem 3.2: the problem should have read ``$\dots$but
$x^{3}+x-4\notin C$'' \uline{not} ``$\dots$but $x^{3}+x-\notin C$''.
\end{itemize}
\begin{defn}
\textbf{[]Exercise \ref{exc:constant_polys_are_ideals?} }Let $\mathbb{R}[x]$,
the set of all polynomials with real coefficients. Let $C$ be the
subset of $\mathbb{R}[x]$ with constant terms that are $0$. For
example, $x^{3}+x\in C$ but $x^{3}+x-4\notin C$. Prove that $C$
is or is not an ideal. 

\begin{itemize}
\item [Hint:]You should multiply a few polynomials and see what happens.
If it is an ideal, it should pass the ideal test (theorem \ref{thm:Ideal-test.-A}).
If it is not, a simple example will suffice. Theorem \ref{thm:Ideal-test.-A}
is an ``if and only if'' statement. So if part 1 or 2 is not met,
then the non-empty subset is not an ideal. Keep in mind that for part
2, the $r\in R$ means any element in the ring. For this problem,
that would include polynomial with non-zero constants like $x^{3}+x-4$,
$x-2$, $3$, etc. 
\end{itemize}
\end{defn}

\begin{defn}
\textbf{[]Exercise \ref{exc:I=00003DR-if-I-has-an-inverse} }Let
$R$ be a ring with unity $1_{R}$ and let $I$ be an ideal of $R$.
If $I$ has an element $x$ and $x^{-1}$ such that $x\cdot x^{-1}=x^{-1}\cdot x=1_{R}$\footnote{Such an element $x$ is called and \emph{unit.}},
prove that $I=R$.

\begin{itemize}
\item [Hint:]We already know that $I\subseteq R$ so we only need to show
that $R\subseteq I$. Key here is that $1_{R}$ is in $I$, and that
$r=1\cdot r$.
\end{itemize}
\end{defn}

\end{document}
